\section{Progettazione Logica}

\subsection{Tabella dei volumi}

Per valutare i costi delle operazioni, per effettuare una corretta analisi delle ridondanze 
e quindi per effettuare una corretta ristrutturazione dello schema E-R, 
si rende necessaria una tabella dei volumi in cui per ogni oggetto stimiamo 
la sua quantità media presente nella base di dati a metà del suo ciclo di vita operativo.

\begin{table}[h]
    \centering
    \renewcommand{\arraystretch}{1.2}
    \begin{tabular}{|l|c|r|}
        \hline
        \textbf{Concetto} & \textbf{Tipo} & \textbf{Volume} \\
        \hline
        DIPARTIMENTO & E & 3 \\
        \hline
        CORSO DI LAUREA & E & 21 \\
        \hline
        DOCENTE & E & 320 \\
        \hline
        CORSO & E & 350 \\
        \hline
        EDIZIONE & E & 700 \\
        \hline
        LEZIONE & E & 14.000 \\
        \hline
        STUDENTE & E & 6.500 \\
        \hline
        ESAME & E & 78.000 \\
        \hline
        è affiliato ad & R & 320 \\
        \hline
        è iscritto & R & 6.500 \\
        \hline
        nel piano di studio & R & 130.000 \\
        \hline
        sostiene & R & 78.000 \\
        \hline
        prevede & R & 78.000 \\
        \hline
        prerequisito & R & 100 \\
        \hline
        può insegnare & R & 960 \\
        \hline
        è erogato in & R & 700 \\
        \hline
        tiene & R & 850 \\
        \hline
        si compone di & R & 14.000 \\
        \hline
    \end{tabular}
    \caption{Tabella dei volumi stimati per il Polo Scientifico}
    \label{tab:volumi}
\end{table}

\noindent 
La stima dei volumi (Tabella \ref{tab:volumi}) è stata ricavata da un'analisi quantitativa legata a una realtà accademica concreta, 
prendendo come modello di riferimento i tre dipartimenti del Polo Scientifico dell'Università degli Studi di Udine (\texttt{DMIF}, \texttt{DPIA}, \texttt{DI4A}). 
I dati di riferimento sono stati estrapolati e proporzionati partendo dagli Open Data ufficiali del Ministero dell'Università e della Ricerca (portale USTAT) 
relativi all'anno 2024 e dall'offerta formativa di ateneo per l'anno accademico 2023/2024. 
Sapendo che l'università ospita storicamente circa 16.000 iscritti complessivi e che l'area STEM (Informatica, Ingegneria, Scienze Agrarie) 
assorbe tipicamente il 40\% del totale degli iscritti, si è stimato un volume di circa 6.500 studenti attivi. 
Allo stesso modo, l'organico di docenti strutturati afferenti a questi tre dipartimenti ammonta a circa 320 unità, 
i quali gestiscono un'offerta di 21 corsi di laurea (tra triennali e magistrali) e 350 insegnamenti unici. \\

\noindent
Partendo da queste basi fattuali, i volumi delle relazioni e delle entità intermedie sono stati ricavati matematicamente per un sistema a metà del suo ciclo di vita:

\begin{itemize}
    \item \textbf{EDIZIONE e LEZIONE}: Mantenendo uno storico operativo di due anni accademici, i 350 corsi generano 700 istanze per l'entità EDIZIONE e per la relazione \textit{è erogato in}. Assumendo una media di 20 appuntamenti per edizione, si ricavano $700 \cdot 20 = 14.000$ istanze per l'entità LEZIONE e per la relazione \textit{si compone di}.
    \item \textbf{Piano di studio}: Si stima che uno studente inserisca in media 20 corsi nel proprio piano di studi durante la sua carriera. Considerando 6.500 studenti attivi, la relazione \textit{nel piano di studio} conta $6.500 \cdot 20 = 130.000$ tuple.
    \item \textbf{ESAME, sostiene e prevede}: Considerando l'eterogeneità della popolazione studentesca, costituita da neo-immatricolati, studenti in pari e fuori corso, si stima una media ponderata di 12 esami superati per studente. Di conseguenza, l'entità ESAME e le relazioni \textit{sostiene} e \textit{prevede} generano 6.500 $\cdot$ 12 = 78.000 istanze.
    \item \textbf{può insegnare e tiene}: Ipotizzando che ogni docente sia abilitato a insegnare mediamente 3 corsi afferenti al proprio settore disciplinare, la relazione \textit{può insegnare} genera $320 \cdot 3 = 960$ tuple. Per la relazione \textit{tiene}, includendo le co-docenze, si stimano 850 assegnazioni sulle 700 edizioni attive.
    \item \textbf{prerequisito}: Si ipotizza che all'interno di ciascuno dei 21 corsi di laurea vi siano in media 4 o 5 esami propedeutici per poter essere sostenuti, generando una stima di 100 istanze per la relazione ricorsiva.
\end{itemize}

\subsection{Analisi delle ridondanze}

In questa sezione viene valutata l'opportunità di introdurre una ridondanza controllata all'interno dello schema logico. 
Si evidenzia che lo schema concettuale Entità-Relazioni (E-R) precedentemente definito è rigorosamente privo di ridondanze, 
tutte le entità in esso presenti contengono esclusivamente attributi atomici e non derivabili. 
In questa fase di progettazione logica, tuttavia, si analizza l'impatto prestazionale derivante dall'introduzione intenzionale di un attributo ridondante, 
denominato \textit{NumeroIscritti}, all'interno di \textit{CORSO}. \\ 

\noindent
Tale scelta comporta una denormalizzazione del modello, poiché l'informazione legata al numero di studenti 
che hanno scelto un determinato insegnamento è già interamente ricavabile in modo dinamico contando le tuple della relazione \textit{nel piano di studio}. 
L'obiettivo dell'analisi è stabilire se il risparmio computazionale in lettura giustifichi il sovraccarico in scrittura e la potenziale complessità di gestione delle anomalie di modifica. \\

\noindent 
I due task principali selezionati per valutare l'impatto di questa modifica sono:
\begin{itemize}
    \item \textbf{Task 1 (Calcolo iscritti a un corso specifico)}: Dato un corso $X$, restituire il numero totale di studenti che lo hanno inserito nel proprio piano di studi (corrisponde all'OPERAZIONE 1).
    \item \textbf{Task 2 (Inserimento di un corso nel piano di studi)}: Dato un corso $X$, inserirlo nel piano di studi di un determinato studente (corrisponde all'OPERAZIONE 4).
\end{itemize}

\noindent 
In base ai parametri dimensionali stabiliti, si stima che il Task 1 venga eseguito con una frequenza di 20.000 volte all'anno da parte di docenti e segreterie per il monitoraggio delle classi. 
Il Task 2 viene eseguito interattivamente dagli studenti principalmente a ridosso dell'inizio dei semestri, registrando circa 32.500 modifiche annue ($6.500 \text{ studenti} \cdot 5 \text{ corsi/anno}$).

\subsubsection{Accessi in assenza della ridondanza}

\paragraph{Task 1:} In assenza dell'attributo ridondante, per contare gli iscritti a un corso $X$, il DBMS deve scansionare la relazione \textit{nel piano di studio}. Avendo stimato un volume totale di 130.000 tuple distribuite su 350 corsi, l'operazione richiede mediamente l'accesso in lettura a $130.000 / 350 \approx 371$ record per ogni interrogazione.
$$371 \text{ letture} \cdot 20.000 \text{ volte/anno} = 7.420.000 \text{ accessi in lettura all'anno}$$

\paragraph{Task 2:} L'inserimento del corso richiede esclusivamente la scrittura di una nuova tupla all'interno della relazione \textit{nel piano di studio}, senza impattare altre strutture.
$$1 \text{ scrittura} \cdot 32.500 \text{ volte/anno} = 32.500 \text{ accessi in scrittura all'anno}$$

\subsubsection{Accessi in presenza della ridondanza}

\paragraph{Task 1:} Introducendo l'attributo ridondante, il sistema non deve più navigare nella relazione intermedia ed è sufficiente accedere alla singola istanza dell'entità \textit{CORSO} corrispondente all'insegnamento $X$ e leggerne direttamente il valore memorizzato.
$$1 \text{ lettura} \cdot 20.000 \text{ volte/anno} = 20.000 \text{ accessi in lettura all'anno}$$

\paragraph{Task 2:} In presenza di ridondanza, l'operazione di scrittura si complica. Oltre all'inserimento del record nella relazione \textit{nel piano di studio} (1 scrittura), il sistema deve aggiornare il contatore nell'entità \textit{CORSO}, richiedendo una lettura preliminare del valore corrente e la successiva riscrittura del dato incrementato (1 lettura e 1 scrittura). Ogni esecuzione comporta quindi 1 lettura e 2 scritture complessive.
$$32.500 \text{ letture/anno e } 65.000 \text{ scritture/anno}$$

\subsubsection{Decisione sulla ridondanza}

Al fine di confrontare in modo omogeneo i due scenari, le operazioni di modifica vengono pesate con un costo computazionale doppio rispetto alle operazioni di consultazione ($1w = 2r$). 
I costi complessivi annui espressi in accessi equivalenti in lettura sono riassunti nella seguente tabella:

\begin{table}[h]
    \centering
    \renewcommand{\arraystretch}{1.2}
    \begin{tabular}{|l|r|r|}
        \hline
        \textbf{Task} & \textbf{Con ridondanza} & \textbf{Senza ridondanza} \\
        \hline
        Task 1 & 20.000 & 7.420.000 \\
        \hline
        Task 2 & 162.500 & 65.000 \\
        \hline
        \textbf{Totale [Accessi equivalenti]} & \textbf{182.500} & \textbf{7.485.000} \\
        \hline
    \end{tabular}
    \caption{Bilancio complessivo degli accessi}
\end{table}

\noindent
Il bilancio dei costi evidenzia un divario macroscopico, a fronte di un incremento controllato dei costi di scrittura per il Task 2 (+97.500 accessi equivalenti), 
il Task 1 beneficia di un abbattimento drastico dei tempi di risposta, consentendo di risparmiare oltre 7,3 milioni di accessi annui sul disco. \\

\noindent 
Dal punto di vista progettuale, si decide pertanto di accettare il compromesso della denormalizzazione, modificando formalmente lo schema logico rispetto a quello concettuale originario. 
L'attributo ridondante \textit{NumeroIscritti} viene ufficialmente introdotto all'interno dell'entità \textit{CORSO}, 
delegando a opportuni vincoli di integrità a livello logico (tramite trigger in fase di implementazione) il compito di mantenere aggiornato il contatore ed evitare qualsiasi rischio di inconsistenza della base di dati.

\subsubsection{Analisi dei cicli}

L'analisi delle ridondanze si è estesa alla verifica strutturale dello schema concettuale per l'individuazione di eventuali cicli. 
All'interno del diagramma E-R sono stati individuati due percorsi ciclici principali, i quali sono stati analizzati per determinare la presenza di derivabilità.

\paragraph{Ciclo 1:}
Il primo ciclo si instaura tra le entità \textit{STUDENTE}, \textit{CORSO} ed \textit{ESAME}. È possibile collegare uno studente a un corso tramite due percorsi distinti:

\begin{enumerate}
    \item L'associazione diretta tramite la relazione \textit{nel piano di studio}.
    \item La navigazione indiretta tramite le relazioni \textit{sostiene} e \textit{prevede}.
\end{enumerate}

\noindent
Il ciclo \textbf{non rappresenta una ridondanza}, perché i due percorsi descrivono informazioni diverse.
La relazione relativa ai corsi seguiti dallo studente ha uno scopo programmatorio, in quanto indica le attività che lo studente ha scelto nel proprio percorso di studi.
La relazione legata agli esami sostenuti, invece, ha uno scopo storico, poiché registra le prove realmente effettuate e concluse.
Di conseguenza, eliminare una delle due relazioni comporterebbe la perdita di informazioni importanti.
La coerenza tra i due percorsi, ad esempio il fatto che uno studente possa sostenere solo esami relativi a corsi presenti nel proprio piano di studi, 
viene garantita tramite un vincolo di integrità implementato con un trigger, e non attraverso modifiche alla struttura dello schema.

\paragraph{Ciclo 2:}
Il secondo ciclo si instaura tra le entità \textit{DOCENTE}, \textit{CORSO} ed \textit{EDIZIONE}. Un docente risulta collegato a un corso tramite:

\begin{enumerate}
    \item L'associazione diretta tramite la relazione \textit{può insegnare}.
    \item La navigazione indiretta tramite le relazioni \textit{tiene} ed \textit{è erogato in}.
\end{enumerate}

\noindent
Anche in questo caso il ciclo \textbf{non costituisce ridondanza}. 
La relazione \textit{può insegnare} modella un'abilitazione di base, indicando le materie afferenti al settore scientifico-disciplinare del docente, 
indipendentemente dall'effettiva erogazione. Il percorso indiretto modella invece l'incarico reale assegnato in un preciso anno accademico. 
Essendo l'abilitazione (statica) e l'assegnazione (dinamica) due concetti temporalmente e funzionalmente separati, si decide di mantenere intatto lo schema E-R.

\subsection{Schema Logico}
Di seguito si riporta lo schema logico risultante. Per convenzione, gli attributi che compongono la \underline{chiave primaria} sono sottolineati, mentre le chiavi esterne sono contrassegnate da un asterisco (*).

\begin{itemize}
    \item \textbf{DIPARTIMENTO} (\underline{CodiceDipartimento}, Nome)
    \item \textbf{CORSO\_DI\_LAUREA} (\underline{CodiceCDL}, Nome)
    \item \textbf{DOCENTE} (\underline{CodiceFiscale}, Nome, Cognome, Titolo, CodiceDipartimento*)
    \item \textbf{STUDENTE} (\underline{Matricola}, Nome, Cognome, Telefono, AA\_Immatricolazione, CodiceCDL*)
    \item \textbf{CORSO} (\underline{CodiceCorso}, Nome, Crediti, Descrizione, NumeroIscritti)
    \item \textbf{EDIZIONE} (\underline{CodiceCorso*}, \underline{AnnoAccademico}, \underline{Periodo})
    \item \textbf{LEZIONE} (\underline{CodiceCorso*}, \underline{AnnoAccademico*}, \underline{Periodo*}, \underline{Data}, \underline{FasciaOraria}, Aula)
    \item \textbf{ESAME} (\underline{Matricola*}, \underline{CodiceCorso*}, \underline{Data}, Punteggio)
    \item \textbf{PIANO\_DI\_STUDIO} (\underline{Matricola*}, \underline{CodiceCorso*})
    \item \textbf{PREREQUISITO} (\underline{CodiceCorsoPropedeutico*}, \underline{CodiceCorsoSuccessivo*})
    \item \textbf{PUO\_INSEGNARE} (\underline{CodiceFiscale*}, \underline{CodiceCorso*})
    \item \textbf{TIENE} (\underline{CodiceFiscale*}, \underline{CodiceCorso*}, \underline{AnnoAccademico*}, \underline{Periodo*})
\end{itemize}

\subsection{Vincoli di Integrità Referenziale}
Per garantire la coerenza della base di dati, si definiscono esplicitamente i seguenti vincoli di integrità referenziale (Foreign Keys) per lo schema logico appena descritto:

\begin{itemize}
    \item \textbf{DOCENTE}.CodiceDipartimento $\rightarrow$ \textbf{DIPARTIMENTO}.CodiceDipartimento
    \item \textbf{STUDENTE}.CodiceCDL $\rightarrow$ \textbf{CORSO\_DI\_LAUREA}.CodiceCDL
    \item \textbf{EDIZIONE}.CodiceCorso $\rightarrow$ \textbf{CORSO}.CodiceCorso
    \item \textbf{LEZIONE}.(CodiceCorso, AnnoAccademico, Periodo) $\rightarrow$ \textbf{EDIZIONE}.(CodiceCorso, AnnoAccademico, Periodo)
    \item \textbf{ESAME}.Matricola $\rightarrow$ \textbf{STUDENTE}.Matricola
    \item \textbf{ESAME}.CodiceCorso $\rightarrow$ \textbf{CORSO}.CodiceCorso
    \item \textbf{PIANO\_DI\_STUDIO}.Matricola $\rightarrow$ \textbf{STUDENTE}.Matricola
    \item \textbf{PIANO\_DI\_STUDIO}.CodiceCorso $\rightarrow$ \textbf{CORSO}.CodiceCorso
    \item \textbf{PREREQUISITO}.CodiceCorsoPropedeutico $\rightarrow$ \textbf{CORSO}.CodiceCorso
    \item \textbf{PREREQUISITO}.CodiceCorsoSuccessivo $\rightarrow$ \textbf{CORSO}.CodiceCorso
    \item \textbf{PUO\_INSEGNARE}.CodiceFiscale $\rightarrow$ \textbf{DOCENTE}.CodiceFiscale
    \item \textbf{PUO\_INSEGNARE}.CodiceCorso $\rightarrow$ \textbf{CORSO}.CodiceCorso
    \item \textbf{TIENE}.CodiceFiscale $\rightarrow$ \textbf{DOCENTE}.CodiceFiscale
    \item \textbf{TIENE}.(CodiceCorso, AnnoAccademico, Periodo) $\rightarrow$ \textbf{EDIZIONE}.(CodiceCorso, AnnoAccademico, Periodo)
\end{itemize}
